% GENERATED FILE. Edit canonical lessons, then run npm run generate:books.
% canonical-source-hash: fnv1a64:5e98abc470a4866e
% canonical-lessons: MW-C50-lugai, MW-C50-parosi, MW-C50-pavno, MW-C50-mastar, MW-C50-khano

\chapter{Woman, Neighbour, Guest, Teacher, To eat}
\label{ch:mw-woman-neighbour-guest-teacher-to-eat}
\hlchaptermodality{50}

\begin{chapteropening}
\textbf{By the end of this chapter:} \emph{I can say a woman, a neighbour, a guest, a teacher and to eat.}

Complete the last lesson of chapter 50: I can say a woman, a neighbour, a guest, a teacher and to eat.
\end{chapteropening}

\section[\texorpdfstring{\mw{लुगाई} (lugāī)}{लुगाई (lugāī)}]{\mw{लुगाई} (lugāī) — a woman}

\label{lesson:MW-C50-lugai}



\begin{quote}
Before the new one: say the Marwadi for a girl, then the Marwadi for a man.
\end{quote}

\subsection*{You'll want to know: \mw{लुगाई}}
\textbf{\mw{लुगाई}} — \emph{lugāī} — \textquotedblleft{}a woman\textquotedblright{}.

\begin{cousinweb}
The origin of \textbf{\mw{लुगाई}} (\emph{lugāī}) is not settled. CDIAL, searched by form and by meaning, has no entry for it.
\end{cousinweb}

\begin{grammarlens}[title={What you've built}]
One more word for this chapter.
\end{grammarlens}

\subsection*{Your turn}
\begin{itemize}
\raggedright
  \item \emph{Say it:} \emph{lugāī}
  \item \emph{Say it:} \emph{lugāī}, once more
  \item \emph{Say it:} say \emph{lugāī}
  \item \emph{Recall:} say the Marwadi for a girl, then the Marwadi for a man, then say \emph{lugāī} again
\end{itemize}

\begin{tcolorbox}[breakable,colback=teal!4,colframe=teal!35!black,title={Before you move on}]
What does \mw{लुगाई} mean? (\textquotedblleft{}A woman\textquotedblright{}.) Say it once more.
\end{tcolorbox}
\section[\texorpdfstring{\mw{पाड़ोसी} (pāṛosī)}{पाड़ोसी (pāṛosī)}]{\mw{पाड़ोसी} (pāṛosī) — a neighbour}

\label{lesson:MW-C50-parosi}



\begin{quote}
Before the new one: say the Marwadi for a man, then the Marwadi for a woman.
\end{quote}

\subsection*{You'll want to know: \mw{पाड़ोसी}}
\textbf{\mw{पाड़ोसी}} — \emph{pāṛosī} — \textquotedblleft{}a neighbour\textquotedblright{}.

\begin{cousinweb}
\textbf{\mw{पाड़ोसी}} (\emph{pāṛosī}) is a tadbhava, a word inherited from Sanskrit and reshaped by sound change. It comes from \emph{prātiveśya}, \textquotedblleft{}neighbouring; neighbour\textquotedblright{}. CDIAL gives no Marwari form in this entry (CDIAL 8936); the match is through Gujarati \emph{pāṛośī} and Sindhi \emph{pāṛosī}. It is built on \emph{prati}, \textquotedblleft{}near\textquotedblright{}, and \emph{veśa}, \textquotedblleft{}house\textquotedblright{}.
\end{cousinweb}

\begin{grammarlens}[title={What you've built}]
One more word for this chapter.
\end{grammarlens}

\subsection*{Your turn}
\begin{itemize}
\raggedright
  \item \emph{Say it:} \emph{pāṛosī}
  \item \emph{Say it:} \emph{pāṛosī}, once more
  \item \emph{Say it:} say \emph{pāṛosī}
  \item \emph{Recall:} say the Marwadi for a man, then the Marwadi for a woman, then say \emph{pāṛosī} again
\end{itemize}

\begin{tcolorbox}[breakable,colback=teal!4,colframe=teal!35!black,title={Before you move on}]
What does \mw{पाड़ोसी} mean? (\textquotedblleft{}A neighbour\textquotedblright{}.) Say it once more.
\end{tcolorbox}
\section[\texorpdfstring{\mw{पावणो} (pāvṇo)}{पावणो (pāvṇo)}]{\mw{पावणो} (pāvṇo) — a guest}

\label{lesson:MW-C50-pavno}



\begin{quote}
Before the new one: say the Marwadi for a woman, then the Marwadi for a neighbour.
\end{quote}

\subsection*{You'll want to know: \mw{पावणो}}
\textbf{\mw{पावणो}} — \emph{pāvṇo} — \textquotedblleft{}a guest\textquotedblright{}.

\begin{cousinweb}
\textbf{\mw{पावणो}} (\emph{pāvṇo}) is a tadbhava, a word inherited from Sanskrit and reshaped by sound change. It comes from \emph{prāhuṇa}, \textquotedblleft{}guest\textquotedblright{}. CDIAL lists Marwari \emph{pāhuṇo} under it (CDIAL 8973). \textbf{\mw{पावणो}} (\emph{pāvṇo}) is the same word said with \emph{v} for \emph{h}.
\end{cousinweb}

\begin{grammarlens}[title={What you've built}]
One more word for this chapter.
\end{grammarlens}

\subsection*{Your turn}
\begin{itemize}
\raggedright
  \item \emph{Say it:} \emph{pāvṇo}
  \item \emph{Say it:} \emph{pāvṇo}, once more
  \item \emph{Say it:} say \emph{pāvṇo}
  \item \emph{Recall:} say the Marwadi for a woman, then the Marwadi for a neighbour, then say \emph{pāvṇo} again
\end{itemize}

\begin{tcolorbox}[breakable,colback=teal!4,colframe=teal!35!black,title={Before you move on}]
What does \mw{पावणो} mean? (\textquotedblleft{}A guest\textquotedblright{}.) Say it once more.
\end{tcolorbox}
\section[\texorpdfstring{\mw{मास्टर} (māsṭar)}{मास्टर (māsṭar)}]{\mw{मास्टर} (māsṭar) — a teacher}

\label{lesson:MW-C50-mastar}



\begin{quote}
Before the new one: say the Marwadi for a neighbour, then the Marwadi for a guest.
\end{quote}

\subsection*{You'll want to know: \mw{मास्टर}}
\textbf{\mw{मास्टर}} — \emph{māsṭar} — \textquotedblleft{}a teacher\textquotedblright{}.

\begin{cousinweb}
\textbf{\mw{मास्टर}} (\emph{māsṭar}) is a loan from English \emph{master}.
\end{cousinweb}

\begin{grammarlens}[title={What you've built}]
One more word for this chapter.
\end{grammarlens}

\subsection*{Your turn}
\begin{itemize}
\raggedright
  \item \emph{Say it:} \emph{māsṭar}
  \item \emph{Say it:} \emph{māsṭar}, once more
  \item \emph{Say it:} say \emph{māsṭar}
  \item \emph{Recall:} say the Marwadi for a neighbour, then the Marwadi for a guest, then say \emph{māsṭar} again
\end{itemize}

\begin{tcolorbox}[breakable,colback=teal!4,colframe=teal!35!black,title={Before you move on}]
What does \mw{मास्टर} mean? (\textquotedblleft{}A teacher\textquotedblright{}.) Say it once more.
\end{tcolorbox}
\section[\texorpdfstring{\mw{खाणो} (khāṇo)}{खाणो (khāṇo)}]{\mw{खाणो} (khāṇo) — to eat}

\label{lesson:MW-C50-khano}



\begin{quote}
Before the new one: say the Marwadi for a guest, then the Marwadi for a teacher.
\end{quote}

\subsection*{You'll want to know: \mw{खाणो}}
\textbf{\mw{खाणो}} — \emph{khāṇo} — \textquotedblleft{}to eat\textquotedblright{}.

\begin{cousinweb}
\textbf{\mw{खाणो}} (\emph{khāṇo}) is a tadbhava, a word inherited from Sanskrit and reshaped by sound change. It comes from \emph{khādati}, \textquotedblleft{}chews, bites; eats\textquotedblright{}. CDIAL lists Marwari \emph{khābo} under it (CDIAL 3865). \emph{khābo} is the same verb with another infinitive ending, \emph{-bo}.
\end{cousinweb}

\begin{grammarlens}[title={What you've built}]
One more word for this chapter.
\end{grammarlens}

\subsection*{Your turn}
\begin{itemize}
\raggedright
  \item \emph{Say it:} \emph{khāṇo}
  \item \emph{Say it:} \emph{khāṇo}, once more
  \item \emph{Say it:} say \emph{khāṇo}
  \item \emph{Recall:} say the Marwadi for a guest, then the Marwadi for a teacher, then say \emph{khāṇo} again
\end{itemize}

\begin{tcolorbox}[breakable,colback=teal!4,colframe=teal!35!black,title={Before you move on}]
What does \mw{खाणो} mean? (\textquotedblleft{}To eat\textquotedblright{}.) Say it once more.
\end{tcolorbox}
